\section{Introduction}
\label{sec:introduction}

Conversational AI companions are increasingly used in social and emotionally
oriented interactions. Prior research reports that users may turn to these
systems in response to loneliness and may perceive them as sources of social
support~\cite{defreitasAICompanionsReduce2024,maplesLonelinessSuicideMitigation2024}.
Studies have also examined how humanlike qualities, trust, and perceived
consciousness shape users' relationships with conversational systems~
\cite{seymourExploringInteractionsTrust2021,
guingrichAscribingConsciousnessArtificial2024}.


Research on social chatbots has also examined
friendship formation, self-disclosure, and
relationship development over time~
\cite{croesFriendMitsuku2021,
skjuveChatbotCompanion2021,
skjuveLongitudinalRelationships2022,
skjuveAIFriend2022}.


Longitudinal and qualitative studies describe relationship formation with
social chatbots, while also documenting variation in how users experience
and sustain these relationships~\cite{croesFriendMitsuku2021,
skjuveLongitudinalRelationships2022,skjuveAIFriend2022}. This motivates
examining specific interaction markers and reported duration rather than
assuming that AI companionship has one consistent social effect.

The gap addressed by this project is operational: its research questions
connect author-level, word-normalized emotional-dependence markers with
reported use duration and a separately defined isolation--avoidance index.
This design makes the proposed measures and comparisons explicit while
keeping them distinct from causal claims.

The present project addresses this need through a content-analysis design
focused on student accounts concerning Replika and Character.AI. It asks
whether the frequency of emotional-dependence markers differs between users
reporting less than one month and more than three months of use (RQ2), and
whether reported duration is positively associated with an isolation--
avoidance index (RQ3). The analysis is observational; any association will
not establish a causal effect. The repository's current numerical dataset is
synthetic and serves only to test the computational pipeline.

Research on romantic relationships with
social chatbots further illustrates the
diversity of meanings users attribute
to human--AI relationships~
\cite{panMouRomanticRelationships2024}.


The main contributions of this work are:
\begin{itemize}
  \item A structured content-analysis framework
        for identifying attachment, fear-of-loss,
        daily-support, and social-avoidance markers
        in AI-companion interaction accounts.
  \item A mathematical formulation of normalized
        emotional-dependence and social-avoidance
        measures at the author level.
  \item A reproducible statistical analysis plan
        combining non-parametric group comparisons,
        effect-size estimation, correlation analysis,
        and adjusted regression.
  \item A synthetic-data validation of the
        computational pipeline, explicitly separated
        from future empirical investigation
        involving authentic student accounts.
\end{itemize}



The remainder of this paper is organized as follows.
Section~\ref{sec:methodology} presents the research
methodology and mathematical formulation.
Section~\ref{sec:results} reports the synthetic
pipeline-validation results.
Section~\ref{sec:discussion} discusses their
interpretation and threats to validity.
Section~\ref{sec:conclusion} summarizes the
contributions and outlines future research.
